%! AP Statistics — Unit 2, Lesson 2.7 — Guided Notes
\documentclass[10pt]{article}
\usepackage{apstats-article}
\usepackage{apstats-boxes}

\begin{document}

\pageheader{Unit 2, Lesson 2.7}{Guided Notes}
\namedateperiod

\vspace{0.12in}

% ── OBJECTIVE ────────────────────────────────────────────────────────────────
\begin{objectivebox}
\small By the end of this lesson, I will be able to\dots\\[4pt]
\writeline\\[1pt]
\writeline\\[1pt]
\writeline
\end{objectivebox}

\vspace{0.1in}

% ── VOCABULARY ───────────────────────────────────────────────────────────────
\begin{vocabbox}
\termblanklong{Residual ($e$)}
\termblanklong{Positive residual}
\termblanklong{Negative residual}
\termblanklong{Residual plot}
\termblanklong{Random scatter (in a residual plot)}
\end{vocabbox}

\vspace{0.1in}

% ── HOOK ─────────────────────────────────────────────────────────────────────
\begin{hookbox}
\small\textbf{Opening Question:} The regression equation for predicting average high temperature
($^\circ$F) from daily hours of sunlight is $\widehat{\text{temp}} = 36 + 3(\text{sunlight})$.
A city gets 10 hours of sunlight and the actual high is 66°F.

\vspace{6pt}
What does the model predict? \blank{3cm} \quad How far off is the model? \blank{3cm}

\vspace{4pt}
Is the actual value above or below the regression line? \blank{3cm} \quad How do you know?\\[6pt]
\writeline\\[1pt]\writeline
\end{hookbox}

\vspace{0.1in}

% ── STEP 1: RESIDUAL FORMULA ─────────────────────────────────────────────────
\begin{notesbox}{Step 1 --- The Residual (DAT-1.E)}
\small

\noindent The \textbf{residual} is the difference between the actual value and the predicted value:
\[
  e = \blank{5cm}
\]

\vspace{6pt}
\begin{tabularx}{\linewidth}{@{} >{\bfseries}p{3.4cm} X @{}}
\rowcolor{sky} \textbf{Sign of residual} & \textbf{What it means} \\
Positive ($e > 0$) & Actual value is \blank{3cm} the regression line (model \blank{3cm} the value). \\[4pt]
Negative ($e < 0$) & Actual value is \blank{3cm} the regression line (model \blank{3cm} the value). \\
\end{tabularx}

\vspace{8pt}
\textbf{Interpret} the residual from the hook ($e = 66 - 66 = 0$):
\vspace{4pt}\\
\writeline

\end{notesbox}

\vspace{0.1in}

% ── STEP 2: COMPUTE RESIDUALS ────────────────────────────────────────────────
\begin{notesbox}{Step 2 --- Computing and Verifying Residuals (DAT-1.E)}
\small

\noindent A researcher records average daily sunlight hours ($x$) and average high temperature
($^\circ$F, $y$) for five cities. The regression equation is
$\widehat{\text{temp}} = 36 + 3(\text{sunlight})$.

\vspace{8pt}
\renewcommand{\arraystretch}{1.9}
\begin{tabularx}{\linewidth}{@{} >{\centering\arraybackslash}p{2.6cm}
                                    >{\centering\arraybackslash}p{2.4cm}
                                    >{\centering\arraybackslash}p{2.6cm}
                                    X @{}}
\rowcolor{sky}
\textbf{Sunlight ($x$, hr)} & \textbf{Temp ($y$, °F)} &
\textbf{Predicted ($\hat{y}$)} & \textbf{Residual ($e = y - \hat{y}$)} \\
\midrule
6  & 52 & \blank{2cm} & \blank{2.5cm} \\
8  & 62 & \blank{2cm} & \blank{2.5cm} \\
10 & 66 & \blank{2cm} & \blank{2.5cm} \\
12 & 74 & \blank{2cm} & \blank{2.5cm} \\
14 & 76 & \blank{2cm} & \blank{2.5cm} \\
\midrule
\multicolumn{3}{r}{\textbf{Sum of residuals}} & \blank{2.5cm} \\
\end{tabularx}

\vspace{8pt}
\textbf{Key fact:} For any least-squares regression line, $\sum e_i = $ \blank{1.5cm}.
This is a useful check --- if your sum is not zero, find the arithmetic error.

\end{notesbox}

\vspace{0.1in}

% ── STEP 3: RESIDUAL PLOT ────────────────────────────────────────────────────
\begin{notesbox}{Step 3 --- Constructing a Residual Plot (DAT-1.E)}
\small

\noindent A \textbf{residual plot} is a scatterplot of \blank{3.5cm} (on the $y$-axis)
vs.\ \blank{4.5cm} (on the $x$-axis).

\vspace{8pt}
\textbf{How to construct a residual plot:}
\begin{enumerate}[leftmargin=1.8em, itemsep=2pt]
  \item Draw new axes: $x$-axis = \blank{5cm}, $y$-axis = \blank{3.5cm}.
  \item Plot each point $(x_i,\; e_i)$ from the residual table above.
  \item Draw a dashed horizontal reference line at \blank{2cm}.
\end{enumerate}

\vspace{8pt}
\textbf{Sketch the residual plot for the sunlight/temperature data here:}

\vspace{0.6in}
\begin{center}
\begin{tikzpicture}
  \draw[->] (0,0) -- (8.5,0) node[right]{\small sunlight (hr)};
  \draw[->] (0,-1.4) -- (0,1.4) node[above]{\small residual (°F)};
  \draw[dashed,gray] (0,0) -- (8.5,0);
  % tick marks on x-axis
  \foreach \x/\lbl in {1.2/6, 2.4/8, 3.6/10, 4.8/12, 6.0/14}
    \draw (\x,0.05)--(\x,-0.05) node[below]{\tiny\lbl};
  % tick marks on y-axis
  \foreach \y/\lbl in {-1.2/-2, -0.6/-1, 0.6/1, 1.2/2}
    \draw (0.05,\y)--(-0.05,\y) node[left]{\tiny\lbl};
\end{tikzpicture}
\end{center}

\end{notesbox}

\vspace{0.1in}

% ── STEP 4: INTERPRET ────────────────────────────────────────────────────────
\begin{notesbox}{Step 4 --- Interpreting a Residual Plot (DAT-1.F)}
\small

\noindent The pattern in a residual plot tells you whether a linear model is appropriate.

\vspace{8pt}
\renewcommand{\arraystretch}{1.7}
\begin{tabularx}{\linewidth}{@{} >{\bfseries}p{3.2cm} X @{}}
\rowcolor{sky} \textbf{Pattern} & \textbf{What it means} \\
Random scatter & \blank{9cm} \\
Curved (U-shape) & \blank{9cm} \\
Fanning (spread grows) & \blank{9cm} \\
\end{tabularx}

\vspace{8pt}
\textbf{AP template for a curved residual plot:}\\[4pt]
``The residual plot shows a \blank{3cm} pattern. This suggests that a linear model may
\blank{9cm}.''

\vspace{8pt}
\textbf{Describe the residual plot you constructed above and conclude about the sunlight/temperature model:}\\[6pt]
\writeline\\[2pt]\writeline\\[2pt]\writeline

\end{notesbox}

\vspace{0.1in}

% ── CONNECTIONS ───────────────────────────────────────────────────────────────
\begin{spiralbox}
\small
\begin{multicols}{2}

\textbf{How this connects:}
\begin{itemize}[itemsep=3pt]
  \item Lesson 2.6: $\hat{y}$ from the LSRL is needed to compute a residual.
  \item Lesson 2.8: The LSRL minimizes $\sum e_i^2$ (sum of \textit{squared} residuals).
  \item Unit 3: $r^2$ (coefficient of determination) measures how much variability in $y$
        is explained by the linear model --- it connects to the residuals.
\end{itemize}

\columnbreak

\textbf{Reflection:}\\
Why do we plot residuals vs.\ $x$ rather than just re-drawing the original scatterplot?
What additional information does the residual plot give us?\\[2pt]
\writeline\\[1pt]\writeline\\[1pt]\writeline

\end{multicols}
\end{spiralbox}

\vspace{0.1in}

\begin{notesbox}{Additional Notes}
\vspace{0pt}
\writeline\\\writeline\\\writeline\\\writeline\\\writeline\\\writeline
\end{notesbox}

\end{document}
